%! Minimizing Loss by Gradient Descent — Unit 8, Lesson 8.3 — Lesson Plan
\documentclass[10pt]{article}
\usepackage{linalg-article}
\usepackage{linalg-boxes}
\usepackage{graphicx}   % -article does NOT load graphicx; needed for thumbnails/screenshots
\graphicspath{{images/}}
\newcommand{\TallMath}[1]{$\displaystyle #1\rule[-1.4em]{0pt}{3.2em}$}
\newcommand{\vv}{\mathbf{v}}
\newcommand{\bb}{\mathbf{b}}
\newcommand{\ww}{\mathbf{w}}
\newcommand{\zero}{\mathbf{0}}
\newcommand{\relu}{\mathrm{ReLU}}
\newcommand{\UnitNumberName}{Unit 8: Learning from Data \quad}
\newcommand{\LessonNumberName}{Lesson 8.3: Minimizing Loss by Gradient Descent}

\begin{document}
\begin{center}
    {\Large\bfseries \CourseName: \SchoolYear \\
    {\small\bfseries \UnitNumberName \LessonNumberName}}
\end{center}

\begin{tcolorbox}[colback=blush, colframe=burgundy, boxrule=0.9pt, arc=2mm,
    left=2mm, right=2mm, top=1.5mm, bottom=1.5mm]
    \textbf{Primary Objective:} Students can read a \textbf{loss} $L(w)$ as a bowl whose lowest point is the
    best weight, and use the \textbf{slope} to tell which way is downhill. They can run \textbf{gradient
    descent} --- repeatedly stepping against the slope with the update $w\leftarrow w-s\,L'(w)$ --- to roll a
    weight to the minimum, and explain the role of the \textbf{learning rate} $s$ (too small crawls, too big
    overshoots). They can lift the same rule to many weights with the \textbf{gradient} $\nabla L$
    ($\ww\leftarrow\ww-s\,\nabla L$) and explain why stepping \emph{against} the gradient, and stopping where it
    is $0$, finds the least-error weights. \emph{(Unit~8 is optional enrichment --- explored, not formally
    assessed.)}
\end{tcolorbox}

\renewcommand{\arraystretch}{1.4}
\begin{skillbox}[Priority Ideas \& Skills]{goldbox}
\begin{tabularx}{\linewidth}{@{}X|X@{}}
\textbf{Priority Ideas \& Skills} & \textbf{Key Understandings (the ``why'')} \\[2pt]
\hline
\vspace{-6pt}
\begin{itemize}[leftmargin=1.1em, nosep, after=\vspace{2pt}]
  \item Read a loss $L(w)$ as a \textbf{bowl}; the bottom (slope $=0$) is the best weight.
  \item Use the \textbf{slope} $L'(w)$ to find downhill; step \emph{against} it.
  \item Run the \textbf{gradient-descent update} $w\leftarrow w-s\,L'(w)$ to the minimum.
  \item Lift the rule to many weights with the \textbf{gradient} $\ww\leftarrow\ww-s\,\nabla L$.
\end{itemize}
&
\vspace{-6pt}
\begin{itemize}[leftmargin=1.1em, nosep, after=\vspace{2pt}]
  \item Learning $=$ making the Unit-4 squared-error loss as small as possible.
  \item The slope points \emph{uphill}, so subtracting it moves toward smaller loss.
  \item The \textbf{learning rate} $s$ sets the step: too small crawls, too big overshoots.
  \item A zero slope means flat ground --- the minimum --- so the descent stops there.
\end{itemize}
\end{tabularx}
\end{skillbox}

\begin{skillbox}[{Vocabulary, Concepts, \& Theorems}]{greenbox}
\renewcommand{\arraystretch}{1.3}
\begin{tabularx}{\linewidth}{@{}l X@{}}
\textbf{Loss $L$} & Total (squared) error between predictions and targets; smaller means a better fit. \\
\textbf{Gradient descent} & Repeatedly step downhill on the loss: $w\leftarrow w-s\,L'(w)$. \\
\textbf{Slope $L'(w)$} & How steeply the loss rises; it points \emph{uphill}, so step the opposite way. \\
\textbf{Learning rate $s$} & The size of each step (step size); controls speed vs.\ overshoot. \\
\textbf{Gradient $\nabla L$} & The slopes for all weights at once, packed into a vector (steepest-uphill direction). \\
\textbf{Minimum} & The bottom of the bowl, where the slope is $0$ --- the best (least-error) weights. \\
\end{tabularx}
\end{skillbox}

\begin{fixedskillbox}[Activate Prior Knowledge \& Spiral Review]{blush}
    The Warm-Up rehearses the three moves gradient descent rests on: \textbf{(1)} the \textbf{squared-error
    loss} $\sum(p_i-t_i)^2$ that learning tries to shrink (Unit~4 \& Lesson~8.0) --- here $[3,5]$ vs.\ $[2,6]$
    gives $2$; \textbf{(2)} reading a \textbf{slope} $L'(w)=2w-8$ to find downhill (negative $\Rightarrow$ go
    right, positive $\Rightarrow$ go left); and \textbf{(3)} the \textbf{update arithmetic}
    $w-s\,L'(w)$ (Unit~1) --- from $w=2$ with $s=0.25$ it lands at $w=3$. Debrief item~2 aloud: the slope points
    \emph{uphill}, so we subtract it to descend.
\end{fixedskillbox}

\begin{skillbox}[Hook]{blush}
    Open from Lesson~8.2's preview: we can now \emph{build} networks --- layers, folds, filters --- but nothing
    yet says \emph{which weights} to use. Every choice gives a \textbf{loss} (the Unit-4 sum of squared errors).
    Ask, \textbf{``If good weights mean small loss, how do we \emph{find} the best ones when there are millions
    of them?''} Set the picture: the loss is a \textbf{bowl} (a valley) over the weights, and we want the
    bottom. Pose the fog metaphor: \emph{standing in a valley in thick fog, feeling only the slope underfoot,
    how do you walk to the lowest point?} Name the payoff question of Lesson~8.3: \textbf{``How can a network
    tune millions of weights just by feeling which way is downhill?''} Preview the answer: read the slope, step
    the opposite way, repeat --- gradient descent.
\end{skillbox}

\begin{skillbox}[Lesson]{blush}
\begin{multicols}{2}
\textbf{1. Learning $=$ minimizing loss.} The loss $L$ scores the weights (Unit-4 squared error). Learning is
one job: make $L$ smallest. For a single weight, $L(w)$ is a \textbf{bowl}; use $L(w)=(w-4)^2$ (minimum at
$w=4$). Real nets have millions of weights, so we cannot eyeball the bottom --- we \emph{walk} there.

\textbf{2. Which way is downhill?} The \textbf{slope} $L'(w)=2w-8$ points \emph{uphill}. At $w=0$ it is $-8$
(negative $\Rightarrow$ go right, toward $4$); at the bottom $w=4$ it is $0$. \emph{We give students the slope
--- no calculus.} Use the two-panel figure (bowl $+$ downhill arrow; descent staircase to the minimum).

\columnbreak

\textbf{3. Gradient descent.} The rule $w\leftarrow w-s\,L'(w)$ with learning rate $s=0.25$, from $w=0$:
$0\to2\to3\to3.5\to\cdots\to4$; the loss falls $16\to4\to1\to0.25$. Stop when the slope is $0$ (the step
changes nothing).

\textbf{4. Many weights $+$ why it matters.} The \textbf{gradient} $\nabla L$ is the slopes for all weights at
once; step against it, $\ww\leftarrow\ww-s\,\nabla L$ --- same rule. Learning rate: too small crawls, too big
overshoots. \textbf{Punchline:} a network is Unit~1 (layers) $+$ ReLU (folds) $+$ Unit~4 (loss) $+$ gradient
descent --- no new machinery. \textbf{Road ahead:} 8.4 mean, variance, covariance.
\end{multicols}
\end{skillbox}

\begin{skillbox}[Explicit Instruction: One Gradient-Descent Step]{blush}
\begin{tabularx}{\linewidth}{@{}p{0.46\linewidth} X@{}}
\textbf{Steps (feel \& step).}
\begin{enumerate}[leftmargin=1.3em, nosep, after=\vspace{2pt}]
  \item Evaluate the slope $L'(w)$ at the current weight.
  \item Multiply by the learning rate: $s\,L'(w)$.
  \item \emph{Subtract} it: $w\leftarrow w-s\,L'(w)$ (step against the slope $=$ downhill).
  \item Repeat until the slope is $\approx 0$ (the bottom of the bowl).
\end{enumerate}
&
\textbf{Worked example} ($L(w)=(w-4)^2$, $L'(w)=2w-8$, $s=0.25$). From $w=0$: $L'(0)=-8$, so
$w=0-0.25(-8)=2$; then $L'(2)=-4\Rightarrow w=3$; then $w=3.5$, marching to $w=4$. \textbf{Common slip:}
dropping the minus sign (stepping \emph{with} the slope, going uphill), or multiplying $s\,L'(w)$
\emph{after} subtracting.
\end{tabularx}
\end{skillbox}

\begin{skillbox}[Active Monitoring]{redbox}
Circulate during the notes practice and Tier work. Watch for and correct:
\begin{itemize}[leftmargin=1.2em, nosep]
  \item dropping the minus sign in $w-s\,L'(w)$ (stepping \emph{with} the slope $=$ uphill --- loss grows);
  \item sign errors when evaluating the slope $2w-8$;
  \item multiplying $s\,L'(w)$ \emph{after} subtracting instead of before;
  \item thinking a bigger learning rate is always better (show the overshoot to $w=8$).
\end{itemize}
Cold-call prompts: ``Which way does the slope point?'' ``So which way do we step?'' ``What is the slope at the
bottom?'' ``What does too big a learning rate do?''
\end{skillbox}

\begin{skillbox}[Group Work \& Differentiation]{redbox}
\textbf{Rolling Downhill} activity (tiered; mirrors the notes).
\begin{multicols}{3}
\textbf{Tier R --- Remediate.} Fill the update table from $w=0$ ($0\to2\to3$) and watch the loss fall
($16\to4\to1$).

\columnbreak
\textbf{Tier A --- Approaching.} Descend from the other side ($6\to5\to4.5$); find that the slope is $0$ at the
minimum (the stopping signal); compare learning rates ($s=0.5$ exact vs.\ $s=1$ overshoot to $8$).

\columnbreak
\textbf{Tier E --- Extension.} The \textbf{gradient} of two weights: step $(0,0)\to(0.5,1.5)\to(0.75,2.25)$
toward $(1,3)$; justify stepping against the gradient and stopping at slope $0$.
\end{multicols}
\end{skillbox}

\begin{skillbox}[Individual Work \& Assessment]{redbox}
\textbf{Exit Ticket} (independent, no notes): take one gradient-descent step and check the loss dropped; state
the slope at the minimum and why the descent stops (plus the name ``learning rate''); and a
\textbf{conceptual/justification check} --- in one sentence, why we step \emph{against} the slope. Sort into
``got it / arithmetic slip / concept slip'' to decide whether 8.4 opens with a quick ``feel the slope, step
downhill, stop when flat'' recap. \emph{Enrichment unit --- use as formative feedback, not a graded
assessment.}
\end{skillbox}

\begin{skillbox}[Reinforcement \& Extension]{goldbox}
\textbf{Homework:} take single steps from both sides (toward $w=4$); confirm the loss falls; tune the learning
rate ($s=0.5$ exact vs.\ $s=1$ overshoot); step the gradient of two weights toward $(2,5)$; and short
justifications (step against the slope, slope $=0$ at the best weights, why huge nets must iterate).
\textbf{Extension:} a too-small learning rate (safe but slow) and the \emph{whole-unit synthesis} --- layers
$+$ ReLU $+$ filters $+$ gradient descent $=$ a network that learns. \textbf{Preview:} Lesson~8.4, \emph{Mean,
Variance, and Covariance} --- the shape of the data itself, the last piece of learning from data.
\end{skillbox}
\end{document}
